Work on EV charging demand forecasting has moved in roughly three
waves: classical statistical models in the 2010s, deep recurrent
networks around 2018--2022, and transformer / foundation models from
2023 onwards. We summarise each wave below and place our work inside
the landscape.

\subsection{Classical Statistical Models}

The first generation of papers attacked EV demand the same way the
Box-Jenkins literature attacks any seasonal series: ARMA, ARIMA and
SARIMA, with parameter selection by AIC and residual diagnostics by
Ljung-Box. Khan et al.~\cite{khan2024sarima} report results on the
same Boulder dataset we use; they find that SARIMA consistently beats
plain ARMA / ARIMA because the weekly seasonality cannot be captured
without an explicit seasonal term. Their best MAPE is 0.12 for
SARIMA versus 0.23 for ARIMA. Facebook Prophet has also been used as
an off-the-shelf alternative that handles holidays and seasonality
without manual tuning, although on this kind of series a properly
specified SARIMA tends to match it. The general lesson from this
literature is that the weekly cycle in EV demand is strong enough
that any model that ignores it (random walk, drift, plain ARIMA) ends
up worse than a model that includes the seasonal lag at $s = 7$.

\subsection{Deep Learning: LSTM and Hybrids}

The second wave moved to recurrent networks. LSTM is the workhorse
here: Yi et al.~\cite{yi2022lstm} train LSTM on utility-scale EV
charging data from Utah and Los Angeles, with day-of-week and
temperature as additional inputs, and show that the multivariate
variant clearly beats the univariate one. The recent systematic
review by Alaraj et al.~\cite{alaraj2025survey} confirms that LSTM
(plus its variants, GRU and CNN-LSTM hybrids) is the deep learning
default for this task, and that attention-augmented variants tend to
improve short-horizon accuracy. For datasets of our size
($\sim$2\,000 daily observations), the survey notes that the gains
from heavy hybrids over a well-tuned single LSTM are usually marginal,
which we confirm experimentally in Section~\ref{sec:results}.

\subsection{Foundation Models for Time Series}

Two recent foundation models for time series have appeared since the
proposal was written, and we benchmark both. Google's
TimesFM~\cite{das2024timesfm} is a decoder-only transformer with
patch-based input embedding, pre-trained on a large corpus of public
time series; it ships in two sizes (200M and 500M parameters) and is
designed to be used zero-shot, with the user providing a context
window and reading off the forecast. Amazon's
Chronos~\cite{ansari2024chronos} takes a different approach: it
tokenises the time series via uniform quantisation and runs a T5
encoder-decoder over the resulting token sequence, sampling
probabilistic forecasts (we take the median of 20 samples). Both
models are interesting because they offer competitive baselines
without any training on the target series. Our results
(Section~\ref{sec:results}) suggest that TimesFM is a stronger choice
for this dataset and that bigger Chronos is, surprisingly, worse than
its small sibling.

\subsection{Transformers Trained on EV Data}

Three peer-reviewed papers train transformer-based models directly on
the same Boulder dataset we use. Koohfar et
al.~\cite{koohfar2023transformer} compare a vanilla transformer
against ARIMA, SARIMA, RNN and LSTM on 25 aggregated Boulder stations,
with horizons of 30, 60 and 90 days; they report that the transformer
beats ARIMA by 62--84\% in normalised RMSE.
Kyriakopoulos and Theodoridis~\cite{kyriakopoulos2025comparative}
extend the comparison to four cities (including Boulder) and find that
the transformer dominates at the station level for daily 1--5 day
horizons. Wu et al.~\cite{wu2025catformer} propose a Context-Aware
Temporal Transformer (CAT-Former) and benchmark it against six other
architectures (LSTM, BiLSTM, CNN-LSTM, CNN-BiLSTM, vanilla
transformer, hybrid transformer) on three Boulder stations,
forecasting one hour and one day ahead. CAT-Former is the strongest
model in their study, with MAE in the normalised $[0,1]$ range of
$0.71$--$0.91$ for one-day-ahead forecasts on three different
stations.

Importantly, none of these papers report numbers in the same units or
on the same station selection that we use. Their MAEs in normalised
space are not directly comparable to our raw kWh values. We come back
to this point in Section~\ref{sec:conclusions}; for now we simply note
that the qualitative ordering they all observe (transformer $>$ LSTM
$>$ classical) matches what we find.

\subsection{PatchTST and Our Position}

The architecture we implement in Section~\ref{sec:models} is
PatchTST, introduced by Nie et al.~\cite{nie2023patchtst}. The key
idea is to slice the input series into non-overlapping patches and
treat each patch as a transformer token; this reduces the sequence
length seen by the attention layer (compared to per-step tokens) and
allows the model to learn patch-level patterns. PatchTST is normally
combined with RevIN (Reversible Instance Normalisation,
Kim et al.~\cite{kim2022revin}) to handle the distribution shift that
appears in long-horizon forecasting. The original paper reports state
of the art on several long-horizon time series benchmarks (ETT,
Weather, Electricity), although it has not, to our knowledge, been
applied to EV charging data before.

\begin{table}[!htbp]
\centering
\caption{Summary of related work on EV charging demand forecasting,
focused on studies that use the same Boulder dataset. Numbers are
reported as published; they are not directly comparable to each other
or to ours due to different units, station selections and horizons.}
\label{tab:related}
\setlength{\tabcolsep}{4pt}
\footnotesize
\begin{tabular}{@{}p{2.8cm}p{2.3cm}p{2.0cm}p{1.1cm}p{2.4cm}@{}}
\toprule
Study & Model & Dataset & Metric & Score \\
\midrule
Khan et al.~\cite{khan2024sarima}             & SARIMA       & Boulder      & MAPE & 0.12 \\
Khan et al.~\cite{khan2024sarima}             & ARIMA        & Boulder      & MAPE & 0.23 \\
Koohfar~\cite{koohfar2023transformer}         & Transformer  & Boulder      & RMSE & 0.085--0.112 (n.) \\
Yi et al.~\cite{yi2022lstm}                   & LSTM         & Utah / LA    & MAE  & reported \\
Wu et al.~\cite{wu2025catformer}              & CAT-Former   & Boulder      & MAE  & 0.71--0.91 (n.) \\
Kyriakopoulos~\cite{kyriakopoulos2025comparative} & Transformer & Boulder + 3 & MAE & 0.28--0.45 (n.) \\
\midrule
\textbf{This work} & \textbf{PatchTST + LSTM} & \textbf{Boulder} & \textbf{MAE} & \textbf{25.09 kWh (raw)} \\
\bottomrule
\end{tabular}
\end{table}

The take-away from Table~\ref{tab:related} is that classical SARIMA
and transformer-based models bracket the published literature on this
dataset, with the gap mainly explained by the seasonal component and
the modelling of exogenous variables. Our contribution sits in the
last row of the table: we run PatchTST under the same evaluation
protocol as our SARIMA baselines and find that it beats the
\emph{strongest} non-transformer model in the project (ARIMAX) while
matching the qualitative ordering of the published work.
